\chapter{Mathematical Analysis of the Conventional Method}\label{ch:mathematics}

\section{Contour displacement and differential relations}
\subsection{A partial derivative is not a contour trajectory}
Assume first that $\psi_d$ is continuously differentiable on an open current domain. Changing only excitation while holding both stator currents fixed gives
\begin{equation}
\left.\frac{\partial\psi_d}{\partial i_e}\right|_{i_d,i_q}=L_{de}.
\end{equation}
This experiment generally changes d-flux. A constant-flux displacement asks a different question: how must other currents move to compensate that change?

Parameterize a differentiable path by $\tau$, with currents $i_d(\tau),i_q(\tau),i_e(\tau)$. Applying the chain rule to the composite function gives
\begin{equation}\label{eq:total-derivative}
\frac{\dd\psi_d}{\dd\tau}
=L_{dd}\frac{\dd i_d}{\dd\tau}
+L_{dq}\frac{\dd i_q}{\dd\tau}
+L_{de}\frac{\dd i_e}{\dd\tau}.
\end{equation}
The total derivative includes every input that changes along the path. A partial derivative is one special choice of path, with the other input derivatives zero.

Suppose d-flux is constant, $\psi_d=\lambda$, and $\dd i_e/\dd\tau\ne0$ so that excitation itself can be used as a local path parameter. Divide~\eqref{eq:total-derivative} by $\dd i_e/\dd\tau$:
\begin{align}
0&=L_{dd}\frac{\dd i_d}{\dd i_e}
+L_{dq}\frac{\dd i_q}{\dd i_e}+L_{de},\\
L_{dd}\frac{\dd i_d}{\dd i_e}
&=-L_{de}-L_{dq}\frac{\dd i_q}{\dd i_e}.
\end{align}
If $L_{dd}\ne0$, division and reversal of the sign yield
\begin{equation}\label{eq:contour-total-slope}
\boxed{-\frac{\dd i_d}{\dd i_e}
=\frac{L_{de}}{L_{dd}}
+\frac{L_{dq}}{L_{dd}}\frac{\dd i_q}{\dd i_e}.}
\end{equation}
The first term is excitation compensation by d-current at fixed q-current. The second is additional compensation for whatever change in d-flux is caused by the moving q-current. Cross-saturation enters through $L_{dq}$; a moving q-coordinate alone is not enough if this derivative vanishes.

At fixed $i_q$, $\dd i_q/\dd i_e=0$, and
\begin{equation}\label{eq:fixed-q-slope}
-\left.\frac{\dd i_d}{\dd i_e}\right|_{\psi_d,i_q}
=\frac{L_{de}}{L_{dd}}.
\end{equation}
The implicit-function theorem supplies a local solution for $i_d$ under the nonzero-$L_{dd}$ assumption. It does not guarantee that the same contour branch extends through the entire measured region. Near a vanishing denominator, a d-current compensation slope can become ill-conditioned or cease to be a suitable parameterization.

\subsection{From a differential slope to a finite secant}
For a single admissible constant-flux path connecting excitation values $e_1,e_2$, integrate the derivative:
\begin{align}
i_d(e_2)-i_d(e_1)
&=\int_{e_1}^{e_2}\frac{\dd i_d}{\dd e}\,\dd e,\\
-\frac{\Delta i_d}{\Delta i_e}
&=\frac{1}{e_2-e_1}\int_{e_1}^{e_2}
\left(\frac{L_{de}}{L_{dd}}
+\frac{L_{dq}}{L_{dd}}\frac{\dd i_q}{\dd e}\right)\dd e.
\label{eq:finite-contour-slope}
\end{align}
A finite quotient is therefore a path average, even when both endpoints really lie on the same contour. It equals one fixed local ratio only under further conditions, for example constant sensitivities and fixed q-current. Applying~\eqref{eq:finite-contour-slope} to arbitrary numerical centroids would assume precisely what must first be established: that their flux values are equal and that an admissible connecting contour path exists.

\section{The displacement of contour averages}
The contour statistic changes the whole crossing population when the excitation group changes. Some crossing coordinates move physically; others appear or disappear because a contour is clipped by the measurement boundary or intersects a different mesh. Thus $\Delta\bar i_q/\Delta\bar i_e$ mixes a geometric displacement with changes of the averaging population.

For a hypothetical differentiable family of centroid coordinates, parameterized by its mean excitation $e$, define the full-field value
\begin{equation}
\chi(e)=\psi_d(\bar d(e),\bar q(e),e).
\end{equation}
The same chain rule now gives
\begin{equation}
\chi'(e)=L_{dd}\bar d'(e)+L_{dq}\bar q'(e)+L_{de}.
\end{equation}
Rearranging shows the missing term in a constant-flux interpretation:
\begin{equation}\label{eq:centroid-slope}
-\bar d'(e)=\frac{L_{de}}{L_{dd}}
+\frac{L_{dq}}{L_{dd}}\bar q'(e)-\frac{\chi'(e)}{L_{dd}}.
\end{equation}
The drift term is only one contribution. If the average leaves the contour and its flux value changes, $\chi'$ contributes as well. The measured diagnostic~\eqref{eq:centroid-departure} evaluates a group interpolant, not an independently measured full three-current field at the centroid. It cannot automatically be substituted for $\chi-\lambda$ in~\eqref{eq:centroid-slope}.

An exact finite-endpoint identity is also instructive. Let $\bm c_0,\bm c_1$ be two centroid coordinates and assume the straight segment $\bm c(t)=\bm c_0+t\Delta\bm c$ lies inside a differentiable full-field domain. Then
\begin{equation}
\Delta\psi_d=\int_0^1
\left(L_{dd}(\bm c(t))\Delta\bar d+
L_{dq}(\bm c(t))\Delta\bar q+
L_{de}(\bm c(t))\Delta\bar e\right)\dd t.
\end{equation}
Writing $\langle L_{xy}\rangle=\int_0^1L_{xy}(\bm c(t))\dd t$, and assuming $\Delta\bar e\ne0$ and $\langle L_{dd}\rangle\ne0$, gives
\begin{equation}\label{eq:finite-centroid}
-\frac{\Delta\bar d}{\Delta\bar e}
=\frac{\langle L_{de}\rangle}{\langle L_{dd}\rangle}
+\frac{\langle L_{dq}\rangle}{\langle L_{dd}\rangle}
\frac{\Delta\bar q}{\Delta\bar e}
-\frac{\Delta\psi_d}{\langle L_{dd}\rangle\Delta\bar e}.
\end{equation}
This identity does not assert that the segment is a contour. The last term accounts for its possibly unequal endpoint fluxes. A ratio of path-averaged derivatives in~\eqref{eq:finite-centroid} and an average of derivative ratios along a contour in~\eqref{eq:finite-contour-slope} are different mathematical objects.

\section{Why a mean coordinate need not lie on the contour}
\subsection{An exact curved-contour example}
Consider the analytic field
\begin{equation}\label{eq:parabolic-example}
\psi_d(d,q,e)=L_d(d+k_0e)+c q^2,
\qquad L_d>0.
\end{equation}
Here $L_d$ is a constant in H, $k_0$ is in A/A, and $c$ is in Vs/A$^2$. The field has an exact scalar reduction with coefficient $k_0$. Its d-flux contours at level $\lambda$ satisfy
\begin{align}
L_d(d+k_0e)&=\lambda-cq^2,\\
d&=\frac{\lambda}{L_d}-k_0e-\frac{c}{L_d}q^2.
\label{eq:parabolic-contour}
\end{align}
The contour is curved in the stator-current plane when $c\ne0$. Average~\eqref{eq:parabolic-contour} over any finite set of exact contour points, including measured variation in $e$:
\begin{equation}\label{eq:mean-parabola}
\bar d=\frac{\lambda}{L_d}-k_0\bar e
-\frac{c}{L_d}\overline{q^2},
\qquad
\overline{q^2}=\frac1M\sum_m q_m^2.
\end{equation}
Now evaluate the original field at the average coordinate:
\begin{align}
\psi_d(\bar d,\bar q,\bar e)
&=L_d\left(\frac{\lambda}{L_d}
-\frac{c}{L_d}\overline{q^2}\right)+c\bar q^2\\
&=\lambda-c(\overline{q^2}-\bar q^2).
\end{align}
With population variance $\operatorname{var}(q)=\overline{q^2}-\bar q^2$, the centroid departure is exactly
\begin{equation}\label{eq:variance-departure}
\boxed{\psi_d(\bar d,\bar q,\bar e)-\lambda=-c\,\operatorname{var}(q).}
\end{equation}
For $c>0$, the centroid lies below the chosen flux level. Even a symmetric crossing population with $\bar q=0$ has a nonzero departure whenever its q-coordinates have nonzero spread. Vanishing mean q-current therefore does not eliminate curvature bias.

For two groups evaluated at the same $\lambda$, subtract~\eqref{eq:mean-parabola}:
\begin{equation}
\Delta\bar d=-k_0\Delta\bar e
-\frac{c}{L_d}\Delta\overline{q^2}.
\end{equation}
Hence the signed displacement ratio is
\begin{equation}\label{eq:variance-factor-bias}
-\frac{\Delta\bar d}{\Delta\bar e}
=k_0+\frac{c}{L_d}
\frac{\Delta\overline{q^2}}{\Delta\bar e}.
\end{equation}
The second moment decomposes as
\begin{equation}
\Delta\overline{q^2}
=\Delta\operatorname{var}(q)+\Delta(\bar q^2).
\end{equation}
Both a change of mean position and a change of spread can alter the estimated displacement. If both groups have zero mean q-current, only the variance difference remains. If their second moments are equal, the signed factor is exactly $k_0$ despite a potentially large, identical centroid departure in both groups. A centroid error is therefore not automatically an equally large factor error. The absolute numerator in the conventional estimator requires a further sign check before~\eqref{eq:variance-factor-bias} can be identified with its individual ratio.

\subsection{A local curvature interpretation beyond the example}
Let $\bm\mu=M^{-1}\sum_m\bm c_m$ and $\bm\delta_m=\bm c_m-\bm\mu$. For a sufficiently smooth scalar field,
\begin{equation}
\psi_d(\bm c_m)=\psi_d(\bm\mu)
+\nabla\psi_d(\bm\mu)^{\mathsf T}\bm\delta_m
+\frac12\bm\delta_m^{\mathsf T}H_d(\bm\mu)\bm\delta_m
+R_m,
\end{equation}
where $H_d$ is the Hessian. Average this expression. The linear term vanishes because $\sum_m\bm\delta_m=0$. If all points have exact flux $\lambda$ and
$\Sigma=M^{-1}\sum_m\bm\delta_m\bm\delta_m^{\mathsf T}$, then
\begin{equation}\label{eq:hessian-centroid}
\psi_d(\bm\mu)-\lambda
=-\frac12\operatorname{tr}(H_d(\bm\mu)\Sigma)-\overline R.
\end{equation}
For bounded third derivatives over a sufficiently small neighborhood, $R_m$ is of third order in $\norm{\bm\delta_m}$. Curvature and the spatial distribution of points act together. A large contour, disconnected pieces or strong variation of curvature can invalidate a local second-order approximation, even though the exact parabolic example remains valid for any spread.

\section{Dependence on triangulation and sampling}
The unweighted average is an average with respect to an empirical crossing measure, not an invariant property of the contour geometry. To see this without any machine model, suppose a crossing set has mean $\bm\mu$ and size $M$. Add one duplicate of an existing crossing $\bm c_r$. Its new mean is
\begin{equation}
\bm\mu'=\frac{M\bm\mu+\bm c_r}{M+1},
\qquad
\bm\mu'-\bm\mu=\frac{\bm c_r-\bm\mu}{M+1}.
\end{equation}
The geometric contour has not changed, but its arithmetic mean changes unless $\bm c_r=\bm\mu$. Refining the mesh in one region can have this same weighting effect by creating more intersections there.

If crossings approach a smooth contour parameterized by arclength $\ell$ with nonuniform asymptotic density $\rho(\ell)$, their arithmetic mean approaches
\begin{equation}\label{eq:density-mean}
\frac{\int \bm c(\ell)\rho(\ell)\,\dd\ell}
{\int\rho(\ell)\,\dd\ell}.
\end{equation}
An arclength centroid would instead set $\rho=1$. Equal geometric contours can therefore yield different arithmetic limits under different sampling densities. Disconnected pieces, vertex multiplicity and truncated contours create additional effective weights.

Delaunay connectivity depends on the metric of the triangulation coordinates. Uniform rotation and scaling preserve its Euclidean geometric criterion in the usual nondegenerate case; independent axis rescaling need not do so. Near co-circular configurations, small coordinate changes can switch a diagonal. The intersection population and filtered edges can then change even when the underlying continuous field does not. The conventional triangulation uses physical stator-current axes; the common reduced table uses range-normalized transformed coordinates. These are different numerical operations.

The long-edge rule limits some sparse connections but can preferentially remove one side of a curved or asymmetric contour. No symmetry completion is applied. Consequently, observed q-drift may combine genuine contour movement and discretization or coverage effects. The available statistics do not uniquely apportion these contributions.

\subsection{Why ratios can be sensitive}
For a differentiable ratio $r=A/B$ near $B\ne0$,
\begin{equation}\label{eq:ratio-uncertainty}
\delta r\simeq\frac{\delta A}{B}-\frac{A}{B^2}\delta B,
\end{equation}
so the first-order variance estimate is
\begin{equation}
\operatorname{var}(r)\simeq
\frac{\sigma_A^2}{B^2}
+\frac{A^2\sigma_B^2}{B^4}
-\frac{2A}{B^3}\operatorname{cov}(A,B).
\end{equation}
Small excitation separations amplify both numerator and denominator uncertainty. Averaging many ratios does not remove a systematic displacement bias shared by them.

Even zero-mean errors in $A$ and $B$ need not yield an unbiased ratio. A second-order expansion around nonzero mean $\mu_B$ gives
\begin{equation}
\mathbb E[A/B]\simeq
\frac{\mu_A}{\mu_B}
-\frac{\operatorname{cov}(A,B)}{\mu_B^2}
+\frac{\mu_A\operatorname{var}(B)}{\mu_B^3}.
\end{equation}
These local expansions assume fluctuations small compared with $|\mu_B|$. They do not justify a Gaussian uncertainty model near a denominator crossing zero. The conventional absolute numerator is also nondifferentiable at zero displacement, so its uncertainty needs separate treatment there.

\section{Conditions for an exact common scalar reduction}
\subsection{The discarded direction}
Let $\Omega\subset\mathbb R^3$ be open and $\ps\in C^1(\Omega)$. For a prescribed constant $k$, define
\begin{equation}
\pi_k(i_d,i_q,i_e)=(i_d+k i_e,i_q).
\end{equation}
An exact reduction requires one function $\bm F$ such that
\begin{equation}\label{eq:exact-factorization}
\ps=\bm F\circ\pi_k.
\end{equation}
All physical points with the same $(s,i_q)$ lie on a \emph{fibre}. Moving on that fibre changes excitation by $\dd i_e$ and changes d-current by $-k\,\dd i_e$, with q-current fixed.

Differentiate the two output components of~\eqref{eq:exact-factorization}:
\begin{align}
L_{dd}&=F_{d,s},&L_{de}&=kF_{d,s},\\
L_{qd}&=F_{q,s},&L_{qe}&=kF_{q,s}.
\end{align}
Both outputs must therefore satisfy
\begin{equation}\label{eq:scalar-exactness}
\boxed{L_{de}=kL_{dd},\qquad L_{qe}=kL_{qd}.}
\end{equation}
Identifying a slope from d-flux alone verifies at most the first relation. The second need not follow from it.

\begin{proposition}\label{prop:scalar-factorization}
For a fixed constant $k$, conditions~\eqref{eq:scalar-exactness} are necessary and locally sufficient for~\eqref{eq:exact-factorization}. They are globally sufficient if every nonempty fibre in $\Omega$ is connected. If fibre components are disconnected, equality of flux between those components is an additional requirement.
\end{proposition}
\begin{proof}
Use the invertible change of coordinates $(i_d,i_q,i_e)=(s-ke,q,e)$ and define $\bm h(s,q,e)=\ps(s-ke,q,e)$. Its derivative at fixed $(s,q)$ is
\begin{equation}\label{eq:fibre-derivative}
\frac{\partial\bm h}{\partial e}
=\begin{bmatrix}L_{de}-kL_{dd}\\L_{qe}-kL_{qd}\end{bmatrix}.
\end{equation}
Necessity follows because a function of $(s,q)$ alone cannot vary with $e$. Conversely, if this derivative is zero on an open neighborhood, $\bm h$ is locally independent of $e$, giving a local $\bm F$. Along a connected fibre the fundamental theorem of calculus makes all these values equal. Defining $\bm F$ by that common value is well posed; local coordinate neighborhoods give its $C^1$ regularity on the open projected domain. Zero derivative on one disconnected interval says nothing about the constant value on another interval, which proves the final qualification.
\end{proof}

Connectedness of the whole current domain does not by itself ensure connected fibres. The proposition also requires the \emph{same} $k$ throughout the domain. Finding some suitable local factor at each point is a weaker statement.

\subsection{The sensitivity-vector geometry}
Define the two output sensitivities to d-current and excitation:
\begin{equation}\label{eq:sensitivity-vectors}
\bm a=\begin{bmatrix}L_{dd}\\L_{qd}\end{bmatrix},
\qquad
\bm b=\begin{bmatrix}L_{de}\\L_{qe}\end{bmatrix}.
\end{equation}
Exact local scalar compensation requires $\bm b=k\bm a$. If $\bm a\ne0$, such a scalar exists exactly when
\begin{equation}
\Delta=L_{dd}L_{qe}-L_{qd}L_{de}=0.
\end{equation}
The determinant is the signed area spanned by these vectors in two-output sensitivity space; its units are H$^2$. If both vectors are nonzero, normalize it:
\begin{equation}\label{eq:eta-angle}
\eta=\frac{|\Delta|}{\norm{\bm a}\norm{\bm b}}=|\sin\theta|.
\end{equation}
This measures directional mismatch, not relative flux error. If $\bm a=0$ and $\bm b\ne0$, no d-shift compensates excitation. If both vanish, every scalar satisfies the local equation and the factor is unidentifiable. A vanishing determinant alone would fail to distinguish these cases.

\section{Local and global identification}
\subsection{Deriving the local least-squares factor}
When $\bm b$ is not parallel to $\bm a$, minimize its unexplained sensitivity:
\begin{align}
E(k)&=\norm{\bm b-k\bm a}^2\\
&=\bm b^{\mathsf T}\bm b-2k\bm a^{\mathsf T}\bm b
+k^2\bm a^{\mathsf T}\bm a.
\end{align}
Differentiating and setting the result to zero gives
\begin{align}
E'(k)&=-2\bm a^{\mathsf T}\bm b+2k\bm a^{\mathsf T}\bm a=0,\\
k_{\rm loc}&=\frac{\bm a^{\mathsf T}\bm b}{\bm a^{\mathsf T}\bm a},
\qquad \bm a\ne0.\label{eq:local-ls}
\end{align}
The positive second derivative $2\norm{\bm a}^2$ establishes the unique minimum under this nonzero-sensitivity assumption. The residual $\bm r=\bm b-k_{\rm loc}\bm a$ satisfies $\bm a^{\mathsf T}\bm r=0$: the remaining excitation response is perpendicular to the available compensation direction.

For two components, direct expansion further gives
\begin{equation}
\min_k\norm{\bm b-k\bm a}^2
=\frac{\Delta^2}{\norm{\bm a}^2}.
\end{equation}
This connects determinant and LS residual without identifying either as a finite flux-reconstruction error.

The separate ratios $L_{de}/L_{dd}$ and $L_{qe}/L_{qd}$ each fit one component exactly when their denominator is nonzero. They can disagree. The joint formula~\eqref{eq:local-ls} avoids division by an individually small cross-inductance, but still becomes poorly conditioned when the whole vector $\bm a$ is small.

\subsection{One constant factor fitted to many derivative samples}
For a fixed derivative population $\mathcal V$, minimize
\begin{equation}
E_{\rm glob}(k)=\sum_{j\in\mathcal V}\norm{\bm b_j-k\bm a_j}^2.
\end{equation}
The same differentiation yields
\begin{equation}\label{eq:global-derivative-ls}
k_{\rm der}=
\frac{\sum_j\bm a_j^{\mathsf T}\bm b_j}
{\sum_j\norm{\bm a_j}^2}
=\frac{\sum_j\norm{\bm a_j}^2k_{{\rm loc},j}}
{\sum_j\norm{\bm a_j}^2}.
\end{equation}
Thus the global derivative factor is a sensitivity-weighted mean of local factors, not their arithmetic mean or median. Changing the retained population or current weighting changes the fitted global value.

The reported relative derivative residual has the form
\begin{equation}\label{eq:relative-derivative-residual}
\rho(k)=
\frac{\left(\sum_j\norm{\bm b_j-k\bm a_j}^2\right)^{1/2}}
{\left(\sum_j\norm{\bm b_j}^2\right)^{1/2}}.
\end{equation}
The raw residual is in H; $\rho$ is dimensionless. Neither is an error in Vs. Along an actual discarded-direction path, however,~\eqref{eq:fibre-derivative} can be integrated:
\begin{equation}
\ps(s-ke_2,q,e_2)-\ps(s-ke_1,q,e_1)
=\int_{e_1}^{e_2}(\bm b-k\bm a)\,\dd e.
\end{equation}
Only with suitable pathwise information could one bound a finite discrepancy, for example by $\int\norm{\bm b-k\bm a}\dd e$. A pooled sample residual does not by itself supply that bound.

\subsection{How empirical derivatives enter the argument}
The derivative estimates used here are local affine regressions, not direct inductance measurements. Normalize each measured current by its observed range, choose the nearest $h$ points, and fit the two flux outputs to an intercept plus three normalized-current offsets. Squared normalized distances determine Gaussian weights. Equivalently, the least-squares design and outputs are multiplied by the square roots of those weights. The fitted slopes are divided by the physical current ranges to recover H.

Neighborhood counts $h=40,80,120$ are available, with 80 used for the primary display. A neighborhood must contain at least two nominal excitation groups. Ill-conditioned affine designs and weak sensitivity vectors are excluded by numerical screens. Such rules protect a fit; they do not establish physical validity of every retained derivative. A comparison across neighborhood counts must use their common retained points, which can differ from the primary 80-neighbor population.

\subsection{Why the different factors need not agree}
The conventional factor~\eqref{eq:k-conventional} uses d-flux contours, arithmetic crossing averages and adjacent-group quotients. The local factor~\eqref{eq:local-ls} fits two derivative components at one point. The global factor~\eqref{eq:global-derivative-ls} pools derivative residuals. A reconstruction-optimal constant instead minimizes finite flux errors after a specified common table has been built. A pointwise alignment factor moves a query on an already fixed table to fit its known target flux.

These objectives differ in physical quantity, weights, population, support and numerical representation. They coincide in special exact models with appropriate sampling; agreement is not a general identity. The complete stator differential matrix, needed later for two-axis compensation and control, is
\begin{equation}\label{eq:differential-matrix}
\Lmat=\begin{bmatrix}L_{dd}&L_{dq}\\L_{qd}&L_{qq}\end{bmatrix}.
\end{equation}
The available stator outputs provide this matrix and $\bm b$; they do not provide an independently identified excitation flux linkage or the remaining excitation-circuit inductances.
